\documentclass[11pt]{article}

\usepackage{graphicx, subcaption, amsfonts, amsmath, amsthm, empheq, setspace, lscape, xcolor}
\usepackage[top=1in, bottom=1in, left=1in, right=1in]{geometry}

% define some commands
% command to box formula
\newcommand*\widefbox[1]{\fbox{\hspace{2em}#1\hspace{2em}}}
\newlength\dlf
\newcommand\alignedbox[2]{
  % Argument #1 = before & if there were no box (lhs)
  % Argument #2 = after & if there were no box (rhs)
  &  % Alignment sign of the line
  {
    \settowidth\dlf{$\displaystyle #1$}  
    % The width of \dlf is the width of the lhs, with a displaystyle font
    \addtolength\dlf{\fboxsep+\fboxrule}  
    % Add to it the distance to the box, and the width of the line of the box
    \hspace{-\dlf}  
    % Move everything dlf units to the left, so that & #1 #2 is aligned under #1 & #2
    \boxed{#1 #2}
    % Put a box around lhs and rhs
  }
}

\newcommand\ER{Erd\H{o}s-R'{e}nyi}
\newcommand{\Forall}{\; \forall \;}
\DeclareMathOperator*{\argmin}{\arg\!\min}

% change captions
\captionsetup{width=0.8\textwidth}
% \captionsetup{labelformat=empty,labelsep=none}

% set paragraph indent length
\setlength\parindent{0pt}

% set folder for imported graphics
\graphicspath{ {../figs/temp/} }

\title{Dimensionality of model manifold}

\begin{document}
\maketitle

\section{Background}

The hypothesis being tested is that, while the model manifold is
indeed two-dimensional as Prof. Gear claims, it is characterized by
disparate scales. This is not unlike the model manifold ribbon
explanation given by Sethna and Transtrum. Thus, the system still
exhibits sloppiness despite retaining its two-dimensional model
manifold. \\

If the system truly had two significant parameters in the region of
interest, we would expect a model manifold that looks something like
that in Fig. (1). If we choose a point on this manifold and look at
its intersection with a growing ball, we would consistently find a
two-dimensional sub-manifold which spans a consistent ratio of
parameters $\frac{\theta_1}{\theta_2}$. However, if instead we had one
sloppy parameter, we would still find a two-dimensional model manifold
but it would would be characterized by disparate lengths along the
important and sloppy parameter directions. This is depicted in
Fig. (2). In this case, as we examine the intersection of the manifold
with a growing ball, the shorter, sloppy direction shrinks in
comparison to the important paramter direction. It increasingly
\textit{appears} one-dimensional. \\

To determine which of these descriptions applies to our current
setting, we can perform DMAPS on samples of the model manifold taken
in growing $B_\delta$ around some fixed point. If the second dimension
is pushed down the spectrum as we examine balls of increasing radius,
it indicates that the characteristic lengths of the two dimensions are
not similar, suggesting sloppiness. \\

\begin{figure}[htbp]
  \centering
  \includegraphics[width=\linewidth]{../model-manifold/non-sloppy-illustration}
  \caption{Illustration of non-sloppy two-dimensional model
    manifold. Two-dimensional structure is retained at all scales.}
\end{figure}

\begin{figure}[htbp]
  \centering
  \includegraphics[width=\linewidth]{../model-manifold/sloppy-illustration}
  \caption{Illustration of sloppy two-dimensional model manifold. As
    the ball radius increases, the shorter dimension appears
    diminished with respect to the long axis.}
\end{figure}

\clearpage
\section{Experimental setup}

We choose as base parameters:

\begin{itemize}
  \item $k_1 = 10^{-1}$
  \item $k_{-1} = 10^3$
  \item $k_2 = 10^3$
\end{itemize}

with $A_0 = 1$, $B_0 = C_0 = 0$. Our model response consists of

\begin{align*}
  f(\theta) = \begin{bmatrix} C(t_1) \\ C(t_2) \\ \vdots \\
    C(t_5) \end{bmatrix}
\end{align*}

where $t_i = i/(2 | \lambda_+ |)$ as recommended in Antonios' latest
notes. This ensures we capture the slow timescale, while damping out
the fast dynamics ($\frac{k_1 k_2}{(k_1 + k_{-1} + k_2)^2} \approx
10^{-4}$). Then we sample parameter space in the region

\begin{itemize}
  \item $\log(k_1) \in [-4, -1]$
  \item $\log(k_{-1}) \in [1, 4]$
  \item $\log(k_2) \in [1, 4]$
\end{itemize}

and record the squared error at each point $\delta(\theta) = \|
f(\theta) - f(\theta^*) \|^2$. We then repeat the following for the
three values of $\delta^* = [10^{-5}, 10^{-3}, 10^{-1}]$:

\begin{enumerate}
  \item Randomly select 4000 points such that $\delta(\theta) <
    \delta^*$
  \item Perform DMAPS on the resulting dataset with $\epsilon =
    \frac{\delta^*}{8}$
  \item Record which two eigenvectors, if any, create a
    two-dimensional embedding of the data
\end{enumerate}

We did uncover two-dimensional embeddings at each $\delta^*$, but the
eigenvectors that parameterized these directions were located
increasingly deep in the eigendecomposition. The following lists which
eigenvectors yielded this embedding for each $\delta^*$

\begin{itemize}
  \item $\delta^* = 10^{-5} \rightarrow v_1$ and $v_{14}$
  \item $\delta^* = 10^{-3} \rightarrow v_1$ and $v_{27}$
  \item $\delta^* = 10^{-1} \rightarrow v_1$ and $v_{36}$
\end{itemize}

\section{Summary}
Thus the eigenvector parameterizing the smaller length of the manifold
was pushed lower and lower as $\delta^*$ increased. The actual
embeddings are presented in Figs. (3)-(5) below. The color denotes
$k_{eff}$, and we see that in cases $k_{eff}$ aligns with the first eigenvector. Overall, I believe
this offers good evidence of sloppiness despite the fact that we do
identify a two-dimensional manifold at all scales.

\begin{figure}[htbp]
  \centering
  \includegraphics[width=\linewidth]{./dmapa14}
  \caption{$\phi_{14}$ vs $\phi_{1}$ for $\delta^* = 10^{-5}$, colored
  by $k_{eff}$}
\end{figure}

\begin{figure}[htbp]
  \centering
  \includegraphics[width=\linewidth]{./dmapb27}
  \caption{$\phi_{27}$ vs $\phi_{1}$ for $\delta^* = 10^{-3}$, colored
  by $k_{eff}$}
\end{figure}

\begin{figure}[htbp]
  \centering
  \includegraphics[width=\linewidth]{./dmapc36}
  \caption{$\phi_{36}$ vs $\phi_{1}$ for $\delta^* = 10^{-1}$, colored
  by $k_{eff}$}
\end{figure}

% \bibliographystyle{plain}
% \bibliography{literature.bib}

\end{document}
